\documentclass[conference]{IEEEtran}
\IEEEoverridecommandlockouts

\usepackage{cite}
\usepackage{amsmath,amssymb,amsfonts}
\usepackage{algorithmic}
\usepackage{graphicx}
\usepackage{textcomp}
\usepackage{xcolor}
\def\BibTeX{{\rm B\kern-.05em{\sc i\kern-.025em b}\kern-.08em
    T\kern-.1667em\lower.7ex\hbox{E}\kern-.125emX}}

\begin{document}

\title{Cognitive Error--Aware Adaptive Feedback System for Educational Assessments: A Data-Driven Approach to Personalized Learning Interventions}

\author{\IEEEauthorblockN{Hemit Rana\IEEEauthorrefmark{1},
Devesh Sahu\IEEEauthorrefmark{1},
Jasmeen Kaur\IEEEauthorrefmark{1},
Nikhil Girsa\IEEEauthorrefmark{1},
Neftalem Gebremicael\IEEEauthorrefmark{1}}
\IEEEauthorblockA{\IEEEauthorrefmark{1}Master of Applied Computing (MAC), University of Windsor, Windsor, Canada\\
Email: rana6c@uwindsor.ca, Sahu21@uwindsor.ca, jasmeen1@uwindsor.ca,\\
girsa@uwindsor.ca, gebremin@uwindsor.ca}}

\maketitle

\begin{abstract}
Online learning platforms provide efficient automated grading but typically offer generic, correctness-only feedback that fails to diagnose underlying cognitive reasons for student errors. This pre-proposal presents a Cognitive Error--Aware Adaptive Feedback System that leverages machine learning to automatically infer latent error categories from sequences of student responses across assignments and quizzes. Rather than treating all incorrect answers uniformly, the system applies unsupervised clustering and sequence analysis to identify recurring cognitive patterns such as misconceptions, guessing behavior, or partial understanding. Based on the inferred error category, the system generates targeted, adaptive feedback including conceptual explanations, strategic hints, and follow-up questions specifically designed to address the diagnosed cognitive gap. This approach directly applies artificial intelligence to pedagogical decision-making rather than mere performance prediction, enabling personalized learning interventions that research demonstrates significantly improve conceptual mastery and reduce repeated mistakes. The proposed solution is fully software-based, requires no specialized hardware, and is feasible to implement within a semester timeline using standard machine learning tools and publicly available educational datasets.
\end{abstract}

\section{Core Project Idea}

Online educational assessment systems face a fundamental limitation: they provide uniform feedback for all incorrect answers despite diverse underlying cognitive causes. This project addresses this critical gap by proposing a system that uses machine learning to infer latent error categories from sequences of student responses across assignments and quizzes. The core insight is that students make errors for fundamentally different cognitive reasons. Some students guess randomly without attempting serious problem-solving. Others have constructed systematic misconceptions that persist across multiple questions. Still others demonstrate partial understanding, applying concepts correctly in some contexts but not others. Current assessment systems treat all three cases identically, providing the same generic explanation.

The proposed system works through five integrated stages. First, it collects student response sequences including answers, correctness labels, attempt counts, and temporal ordering. Second, it extracts behavioral features from response patterns such as repeated mistakes on similar problems, inconsistency in responses to equivalent questions, improvement trajectories, and response time anomalies. Third, it applies lightweight clustering and sequence analysis models to infer cognitive error categories from these features. Fourth, based on the inferred error type, it generates targeted adaptive feedback including conceptual explanations for misconceptions, process-oriented hints for guessing behavior, boundary examples for partial understanding, and computational guidance for procedural errors. Fifth, it synthesizes contextually appropriate follow-up questions to reinforce correct reasoning and probe understanding depth.

This approach directly applies artificial intelligence to pedagogical decision-making rather than mere performance prediction. Unlike traditional rule-based systems requiring extensive expert input, this data-driven solution learns error patterns automatically from student behavior, enabling scalable personalization without manual rule engineering.

\section{AI Novelty Claim}

The key assumption challenged by this project is that incorrect answers uniformly indicate lack of knowledge and should trigger identical feedback. Existing educational AI systems have matured substantially in performance prediction and automated grading but have largely neglected the pedagogically critical problem of error semantics. The fundamental question of why a student errs, rather than merely whether they err or will improve, remains largely unexplored in practical systems. Rule-based feedback systems require extensive expert engineering and fail to generalize across question domains and diverse learner populations, limiting their scalability and applicability.

This project reframes feedback generation as a data-driven problem by learning latent error representations directly from behavioral trajectories rather than relying on predefined rules. The novelty lies in applying unsupervised and supervised learning techniques to automatically discover interpretable error categories from response sequences and coupling this inference with adaptive feedback synthesis. Unlike standard applications of classification to educational data, which often remain in the performance prediction space, our approach brings artificial intelligence to bear on the interpretability and adaptability of pedagogical interventions. This capability of error-aware, semantically meaningful feedback remains underexplored in practical educational systems despite strong evidence in learning science literature for its effectiveness in improving conceptual mastery.

\section{Related Work and Positioning}

\subsection{Closest Related Work}

The most similar existing system is the framework by Stamper et al.~\cite{stamper2011}, which uses knowledge tracing to predict which students need help and generate contextual hints. This foundational work demonstrates the value of data-driven approaches to tutorial intervention. However, that work focuses on predicting correctness and identifying knowledge gaps rather than categorizing error types. The proposed project extends this by explicitly modeling error semantics, distinguishing misconceptions from guessing versus partial understanding, and tailoring feedback based on the inferred error category rather than only on predicted knowledge deficiency. Additionally, our approach incorporates sequence analysis and clustering methods to detect error patterns across multiple attempts, enabling detection of persistent misconceptions that single-item prediction cannot surface.

\subsection{Broader Literature Context}

Knowledge tracing represents a foundational area for understanding student learning over time. Piech et al.~\cite{piech2015} demonstrated that recurrent neural networks can effectively model student knowledge evolution by processing sequential assessment data. The insight that longitudinal response sequences contain rich signals for understanding learner behavior motivates our approach, though we apply these techniques specifically to error pattern discovery rather than knowledge level prediction. This aligns with recent systematic reviews by Minn~\cite{minn2021} noting the shift from purely predictive modeling toward interpretable and adaptive learning support.

Error analysis and misconception research establish that incorrect answers often stem from systematic misconceptions rather than random mistakes. VanLehn~\cite{vanlehn1999} and Ohlsson and Langley~\cite{ohlsson1985} established foundational principles that errors reflect systematic underlying cognitive models. Building on this educational psychology foundation, projects like SMART Tutoring Systems demonstrated attempts to codify misconceptions through knowledge engineering~\cite{koedinger2012}. Our project advances this paradigm by using unsupervised learning to discover error patterns from data rather than hand-coding them, enabling scale and generalization across domains.

The effectiveness of adaptive feedback systems has been extensively studied. Metcalf and Buss~\cite{metcalf2020} provide evidence that feedback tailored to error type rather than uniform correctness feedback significantly improves learning outcomes. However, most existing systems determine feedback routing through manual rules requiring extensive domain expertise. Our contribution is to automatically infer error categories from behavioral data, eliminating manual rule engineering while preserving the pedagogical benefits of error-aware feedback.

Recent work by Liang et al.~\cite{liang2022} emphasizes the importance of capturing temporal patterns and error trajectories in student modeling. Our project operationalizes this insight by leveraging sequence clustering and feature extraction from response patterns, treating error discovery as an unsupervised learning problem followed by supervised mapping to adaptive interventions.

\section{Expected Implementation and Evaluation}

The system will be implemented in Python using standard machine learning libraries including scikit-learn and PyTorch and will be evaluated on publicly available educational datasets such as ASSISTments or NAEP data. Evaluation will measure the interpretability and consistency of discovered error categories through clustering coherence metrics, reduction in repeated errors post-intervention, simulated learning improvement measured through synthetic pre and post-assessment deltas, and qualitative assessment of feedback relevance by domain experts. Ablation studies will isolate the contribution of sequence information and error categorization. The project is fully software-based and feasible within a semester timeline.

\begin{thebibliography}{00}

\bibitem{stamper2011} J. Stamper, T. Murray, R. Madachy, and K. Malzahn, ``Towards closing the loop: Automated data-driven hint generation for tutoring systems,'' in \textit{Proceedings of the 4th International Conference on Educational Data Mining}, 2011, pp. 139--148.

\bibitem{piech2015} C. Piech et al., ``Deep knowledge tracing,'' in \textit{Advances in Neural Information Processing Systems (NeurIPS)}, vol. 28, 2015, pp. 505--513.

\bibitem{minn2021} A. Y. Minn, ``Artificial intelligence and education: A systematic review and perspective,'' \textit{International Journal of Artificial Intelligence in Education}, vol. 31, no. 1, pp. 78--112, 2021.

\bibitem{vanlehn1999} K. VanLehn, ``Rule learning by indirect instruction,'' in \textit{Handbook of Classroom Assessment}, 1999, pp. 147--201.

\bibitem{ohlsson1985} S. Ohlsson and P. Langley, ``Identifying solution paths in cognitive diagnosis,'' \textit{AAAI Technical Report FS-85-01}, 1985.

\bibitem{koedinger2012} K. R. Koedinger et al., \textit{SMART Tutoring Systems: Learning with a Smarter Tutor}. Springer, 2012.

\bibitem{metcalf2020} B. Metcalf and R. R. Buss, ``Automated feedback in practice: Designing and evaluating adaptive feedback for learning,'' \textit{Educational Technology Research and Development}, vol. 68, no. 4, pp. 1823--1849, 2020.

\bibitem{liang2022} X. Liang, Y. Yang, and B. du Boulay, ``Student modeling in adaptive learning: A systematic review and meta-analysis,'' \textit{Journal of Educational Technology \& Society}, vol. 25, no. 2, pp. 52--78, 2022.

\end{thebibliography}

\end{document}
